\section*{Capítulo \ref{ch:g3:shapes} --- Figuras y ángulos rectos}

\begin{solution}{exo:g3:shapes:1}
Típico ángulos rectos: las esquinas del tablero, de una hoja de papel,
de la puerta. Un ángulo no recto: la apertura de unas tijeras,
las manecillas del reloj a la 1 en punto. Cada uno se comprueba con el conjunto.
cuadrado (\cref{met:g3:shapes:check}) --- nunca sólo a simple vista.
\end{solution}

\begin{solution}{exo:g3:shapes:2}
Construcción: un lado de $6$ centímetro, ángulos rectos en ambos extremos con el
escuadra, lados de $4$ cm, cerrar la forma. Marcas: $4$ poco
\omterm{def:g3:shapes:shapes}{cuadrados} en las esquinas, una marca en cada $6$ cm de lado, un doble tic
en cada uno $4$ cm de lado.
\end{solution}

\begin{solution}{exo:g3:shapes:3}
un cuadrado tiene $4$ ángulos rectos y $4$ lados iguales --- todos marcados.
\end{solution}

\begin{solution}{exo:g3:shapes:4}
Construcciones de rejilla; la red garantiza la ángulos rectos, y
contar \omterm{def:g3:shapes:shapes}{cuadrados} garantiza las longitudes.
\end{solution}

\begin{solution}{exo:g3:shapes:5}
``Todo cuadrado es un rectángulo'': \emph{verdadero} --- tiene $4$ ángulos rectos, que es todo lo que pide un rectángulo.

``Todo rectángulo es un cuadrado'': \emph{FALSO} --- a $6 \times 4$
El rectángulo tiene lados desiguales.

"Un triángulo puede tener dos ángulos rectos'': \emph{FALSO} --- los dos
lados extraídos de los dos ángulos rectos sería perpendicular a la
misma línea base, por lo tanto nunca se juntan para cerrar el triángulo.
\end{solution}

\begin{solution}{exo:g3:shapes:6}
$ OA = 4$ centímetros (el \omterm{def:g3:shapes:shapes}{radio}). $ OB = 4$ cm también: \emph{cada} punto de
el círculo está a la misma distancia del centro.
\end{solution}

\begin{solution}{exo:g3:shapes:7}
El rectángulo se corta en tres triángulos: dos triángulos rectángulos
($ AMD $ y $ MBC $, en \omterm{def:g3:shapes:rightangle}{ángulo recto} en $ A $ y $ B $) y el \omterm{def:g3:division:half}{medio}
\omterm{def:g3:shapes:shapes}{triángulo} $ MCD $.
\end{solution}

\begin{solution}{exo:g3:shapes:8}
Con al menos uno \omterm{def:g3:shapes:rightangle}{ángulo recto}: el cuadrado, el rectángulo, la derecha
triángulo y la letra L (sus esquinas interior y exterior). el \omterm{def:g3:shapes:shapes}{circulo}
No tiene ningún rincón.
\end{solution}

\begin{solution}{exo:g3:shapes:9}
Cada lado cruza ``$2$ bien, $1$ up'' (o las versiones giradas
``$1$ bien, $2$ abajo'', etc.): cada lado es la diagonal de una
$2 \times 1$ bloque de \omterm{def:g3:shapes:shapes}{cuadrados}, por lo que los cuatro tienen la misma longitud.
\end{solution}

\begin{solution}{exo:g3:shapes:10}
Un programa correcto: ``1. dibujar un rectángulo $ ABCD $ con $ AB = 3$ centímetros
(abajo) y $ a.\,C. = 2$ centímetro. 2. En la parte superior del costado $[DC]$, dibuja el
cuadrado $ DCEF $ de lado $3$ cm, fuera del rectángulo.'' Cualquier claro,
El programa completo, numerado, es correcto. Pruébelo siguiéndolo.
palabra por palabra.
\end{solution}

\begin{solution}{exo:g3:shapes:11}
$1 \times 1$ cuadrícula: $9$. $2 \times 2$ cuadrícula: $4$.
$3 \times 3$: $1$. Total: $9 + 4 + 1 = 14$ cuadrícula.
\end{solution}
